\ifdefined\gkver\else\def\gkver{student}\fi
\documentclass[\gkver]{gaokaozhenti}
\begin{document}

\chapter{2008年全国理综Ⅱ}
%% paperType: 理综物理部分

\begin{enumerate}
\item
%% number: 14
%% paperName: 2008年全国理综Ⅱ第 14 题 6 分
%% typeId: 2
%% type: 多选题
%% chapter: 气体性质
%% point: 气体的状态参量
%% method: 无
%% score: 6
%% degree: 650
%% duplicateId: 0
%% body:
对一定量的气体，下列说法正确的是\xzanswer{BC}

\fourchoices[answer=BC]
{气体的体积是所有气体分子的体积之和}
{气体分子的热运动越剧烈，气体温度就越高}
{气体对器壁的压强是由大量气体分子对器壁不断碰撞而产生的}
{当气体膨胀时，气体分子之间的势能减小，因而气体的内能减少}

%% answer: BC
%% memo:
\memoanswer{本题原卷及材料中未提供详解，待补充。}

\item
%% number: 15
%% paperName: 2008年全国理综Ⅱ第 15 题 6 分
%% typeId: 1
%% type: 单选题
%% chapter: 光学
%% point: 光的折射
%% method: 无
%% score: 6
%% degree: 650
%% duplicateId: 0
%% body:
一束单色光斜射到厚平板玻璃的一个表面上，经两次折射后从玻璃板另一个表面射出，出射光线相对于入射光线侧移了一段距离。在下列情况下，出射光线侧移距离最大的是\xzanswer{D}

\fourchoices[answer=D]
{红光以 $30^{\circ}$ 的入射角入射}
{红光以 $45^{\circ}$ 的入射角入射}
{紫光以 $30^{\circ}$ 的入射角入射}
{紫光以 $45^{\circ}$ 的入射角入射}

%% answer: D
%% memo:
\memoanswer{本题原卷及材料中未提供详解，待补充。}

\item
%% number: 16
%% paperName: 2008年全国理综Ⅱ第 16 题 6 分
%% typeId: 1
%% type: 单选题
%% chapter: 力物体的平衡
%% point: 物体系平衡
%% method: 整体与隔离
%% score: 6
%% degree: 650
%% duplicateId: 0
%% body:
如图，一固定斜面上两个质量相同的小物块 A 和 B 紧挨着匀速下滑，A 与 B 的接触面光滑。已知 A 与斜面之间的动摩擦因数是 B 与斜面之间动摩擦因数的 2 倍，斜面倾角为 $\alpha$。B 与斜面之间的动摩擦因数是\xzanswer{A}

\onepicture[w=4.5cm]{figs/16.png}

\fourchoices[answer=A]
{$\frac{2}{3}\tan\alpha$}
{$\frac{2}{3}\cot\alpha$}
{$\tan\alpha$}
{$\cot\alpha$}

%% answer: A
%% memo:
\memoanswer{本题原卷及材料中未提供详解，待补充。}

\item
%% number: 17
%% paperName: 2008年全国理综Ⅱ第 17 题 6 分
%% typeId: 2
%% type: 多选题
%% chapter: 机械波
%% point: 波的多解性
%% method: 无
%% score: 6
%% degree: 650
%% duplicateId: 0
%% body:
一列简谐横波沿 $x$ 轴正方向传播，振幅为 $A$。$t=0$ 时，平衡位置在 $x=0$ 处的质元位于 $y=0$ 处，且向 $y$ 轴负方向运动；此时，平衡位置在 $x=0.15\Um$ 处的质元位于 $y=A$ 处。该波的波长可能等于\xzanswer{AC}

\fourchoices[answer=AC]
{$0.60\Um$}
{$0.20\Um$}
{$0.12\Um$}
{$0.086\Um$}

%% answer: AC
%% memo:
\memoanswer{本题原卷及材料中未提供详解，待补充。}

\item
%% number: 18
%% paperName: 2008年全国理综Ⅱ第 18 题 6 分
%% typeId: 1
%% type: 单选题
%% chapter: 功和能
%% point: 整体机械能守恒
%% method: 整体与隔离
%% score: 6
%% degree: 550
%% duplicateId: 0
%% body:
如图，一很长的、不可伸长的柔软轻绳跨过光滑定滑轮，绳两端各系一小球 a 和 b。a 球质量为 $m$，静置于地面；b 球质量为 $3m$，用手托住，高度为 $h$，此时轻绳刚好拉紧。从静止开始释放 b 后，a 可能达到的最大高度为\xzanswer{B}

\onepicture[w=2.8cm]{figs/18.png}

\fourchoices[answer=B]
{$h$}
{$1.5h$}
{$2h$}
{$2.5h$}

%% answer: B
%% memo:
\memoanswer{本题原卷及材料中未提供详解，待补充。}

\item
%% number: 19
%% paperName: 2008年全国理综Ⅱ第 19 题 6 分
%% typeId: 1
%% type: 单选题
%% chapter: 电场
%% point: 带电粒子的平衡
%% method: 无
%% score: 6
%% degree: 450
%% duplicateId: 0
%% body:
一平行板电容器的两个极板水平放置，两极板间有一带电量不变的小油滴，油滴在极板间运动时所受空气阻力的大小与其速率成正比。若两极板间电压为零，经一段时间后，油滴以速率 $v$ 匀速下降：若两极板间的电压为 $U$，经一段时间后，油滴以速率 $v$ 匀速上升。若两极板间电压为 $-U$，油滴做匀速运动时速度的大小、方向将是\xzanswer{C}

\fourchoices[answer=C]
{$2v$、向下}
{$2v$、向上}
{$3v$、向下}
{$3v$、向上}

%% answer: C
%% memo:
\memoanswer{本题原卷及材料中未提供详解，待补充。}

\item
%% number: 20
%% paperName: 2008年全国理综Ⅱ第 20 题 6 分
%% typeId: 2
%% type: 多选题
%% chapter: 原子和原子核
%% point: 核聚变
%% method: 无
%% score: 6
%% degree: 650
%% duplicateId: 0
%% body:
中子和质子结合成氘核时，质量亏损为 $\Delta m$，相应的能量 $\Delta E=\Delta mc^{2}=2.2\UMeV$ 是氘核的结合能。下列说法正确的是\xzanswer{AD}

\fourchoices[answer=AD]
{用能量小于 $2.2\UMeV$ 的光子照射静止氘核时，氘核不能分解为一个质子和一个中子}
{用能量等于 $2.2\UMeV$ 的光子照射静止氘核时，氘核可能分解为一个质子和一个中子，它们的动能之和为零}
{用能量大于 $2.2\UMeV$ 的光子照射静止氘核时，氘核可能分解为一个质子和一个中子，它们的动能之和为零}
{用能量大于 $2.2\UMeV$ 的光子照射静止氘核时，氘核可能分解为一个质子和一个中子，它们的动能之和不为零}

%% answer: AD
%% memo:
\memoanswer{本题原卷及材料中未提供详解，待补充。}

\item
%% number: 21
%% paperName: 2008年全国理综Ⅱ第 21 题 6 分
%% typeId: 1
%% type: 单选题
%% chapter: 电磁感应
%% point: 电磁感应中图象
%% method: 无
%% score: 6
%% degree: 450
%% duplicateId: 0
%% body:
如图，一个边长为 $l$ 的正方形虚线框内有垂直于纸面向里的匀强磁场；一个边长也为 $l$ 的正方形导线框所在平面与磁场方向垂直；虚线框对角线 ab 与导线框的一条边垂直，ba 的延长线平分导线框。在 $t=0$ 时，使导线框从图示位置开始以恒定速度沿 ab 方向移动，直到整个导线框离开磁场区域。以 $i$ 表示导线框中感应电流的强度，取逆时针方向为正。下列表示 $i-t$ 关系的图示中，可能正确的是\xzanswer{C}

\onepicture[w=4.6cm]{figs/21a.png}

\fourchoices[answer=C, ispicture=true, h=1.8cm]
{figs/21_A.png}
{figs/21_B.png}
{figs/21_C.png}
{figs/21_D.png}

%% answer: C
%% memo:
\memoanswer{本题原卷及材料中未提供详解，待补充。}

\item
%% number: 22
%% paperName: 2008年全国理综Ⅱ第 22 题 13 分
%% typeId: 4
%% type: 实验题
%% chapter: 电路
%% point: 电路实验
%% method: 无
%% score: 13
%% degree: 450
%% duplicateId: 0
%% body:
\begin{enumerate}
	\item
	（5 分）某同学用螺旋测微器测量一铜丝的直径，螺旋测微器的示数如图所示，该铜丝的直径为 \tkanswer[2cm]{4.593} $\Umm$。

	\onepicture[w=3.6cm]{figs/22.png}

	\item
	（13 分）右图为一电学实验的实物连线图。该实验可用来测量待测电阻 $R_{x}$ 的阻值（约 $500\UO$）。图中两个电压表量程相同，内阻都很大。实验步骤如下：

	①调节电阻箱，使它的阻值 $R_{0}$ 与待测电阻的阻值接近；将滑动变阻器的滑动头调到最右端。

	②合上开关 S。

	③将滑动变阻器的滑动头向左端滑动，使两个电压表指针都有明显偏转。

	④记下两个电压表 $V_{1}$ 和 $V_{2}$ 的读数 $U_{1}$ 和 $U_{2}$。

	⑤多次改变滑动变阻器滑动头的位置，记下 $V_{1}$ 和 $V_{2}$ 的多组读数 $U_{1}$ 和 $U_{2}$。

	⑥求 $R_{x}$ 的平均值。回答下列问题：

	（Ⅰ）根据实物连线图在虚线框内画出实验的电路图，其中电阻箱的符号为图 b、滑动变阻器的符号为图 c，其余器材用通用的符号表示。

	\threepicture[hA=5.0cm, hB=1.3cm, hC=1.3cm,
		capA=实物连线图, capB=电阻箱符号, capC=滑动变阻器符号, align=right]
	{figs/22a.png}
	{figs/22b.png}
	{figs/22c.png}

	（Ⅱ）不计电压表内阻的影响，用 $U_{1}$、$U_{2}$ 和 $R_{0}$ 表示 $R_{x}$ 的公式为 $R_{x}=$ \tkanswer[3cm]{}。

	（Ⅲ）考虑电压表内阻的影响，用 $U_{1}$、$U_{2}$、$R_{0}$、$V_{1}$ 的内阻 $r_{1}$、$V_{2}$ 内阻 $r_{2}$ 表示 $R_{x}$ 的公式为 $R_{x}=$ \tkanswer[4cm]{}。
\end{enumerate}

%% answer:
\jdanswer{
\begin{enumerate}
	\item
	$4.593\Umm$（$4.592$ 或 $4.594$ 也同样给分）

	\item
	（Ⅰ）实验的电路图如图所示。

	\onepicture[w=6.0cm]{figs/22_an.jpg}

	（Ⅱ）$R_{x}=\frac{U_{2}}{U_{1}}R_{0}$

	（Ⅲ）$R_{x}=\frac{U_{2}R_{0}r_{1}r_{2}}{U_{1}r_{2}(R_{0}+r_{1})-U_{2}R_{0}r_{1}}$
\end{enumerate}
}

%% memo:
\memoanswer{本题原卷及材料中未提供详解，待补充。}

\item
%% number: 23
%% paperName: 2008年全国理综Ⅱ第 23 题 15 分
%% typeId: 6
%% type: 计算题
%% chapter: 运动和力
%% point: 动量守恒定律
%% method: 无
%% score: 15
%% degree: 450
%% duplicateId: 0
%% body:
如图，一质量为 $M$ 的物块静止在桌面边缘，桌面离水平地面的高度为 $h$。一质量为 $m$ 的子弹以水平速度 $v_{0}$ 射入物块后，以水平速度 $\frac{1}{2}v_{0}$ 射出。重力加速度为 $g$。求：

\begin{enumerate}
	\item
	此过程中系统损失的机械能；
	\item
	此后物块落地点离桌面边缘的水平距离。
\end{enumerate}

\onepicture[w=4.5cm, align=right]{figs/23.png}

%% answer:
\jdanswer{
\begin{enumerate}
	\item
	$\Delta E=\frac{1}{8}\left(3-\frac{m}{M}\right)mv_{0}^{2}$
	\item
	$s=\frac{mv_{0}}{M}\sqrt{\frac{h}{2g}}$
\end{enumerate}
}

%% memo:
\memoanswer{
（1）设子弹穿过物块后物块的速度为 $V$，由动量守恒得 $mv_{0}=m\cdot\frac{v_{0}}{2}+MV$，解得 $V=\frac{mv_{0}}{2M}$。

系统的机械能损失为 $\Delta E=\frac{1}{2}mv_{0}^{2}-\frac{1}{2}m\left(\frac{v_{0}}{2}\right)^{2}-\frac{1}{2}MV^{2}$，代入得 $\Delta E=\frac{1}{8}\left(3-\frac{m}{M}\right)mv_{0}^{2}$。

（2）设物块下落到地面所需时间为 $t$，落地点距桌边缘的水平距离为 $s$，则 $h=\frac{1}{2}gt^{2}$，$s=Vt$，解得 $s=\frac{mv_{0}}{M}\sqrt{\frac{h}{2g}}$。
}

\item
%% number: 24
%% paperName: 2008年全国理综Ⅱ第 24 题 19 分
%% typeId: 6
%% type: 计算题
%% chapter: 电磁感应
%% point: 电磁感应能量分析
%% method: 无
%% score: 19
%% degree: 650
%% duplicateId: 0
%% body:
如图，一直导体棒质量为 $m$、长为 $l$、电阻为 $r$，其两端放在位于水平面内间距也为 $l$ 的光滑平行导轨上，并与之密接；棒左侧两导轨之间连接一可控制的负载电阻（图中未画出）；导轨置于匀强磁场中，磁场的磁感应强度大小为 $B$，方向垂直于导轨所在平面。开始时，给导体棒一个平行于导轨的初速度 $v_{0}$，在棒的运动速度由 $v_{0}$ 减小至 $v_{1}$ 的过程中，通过控制负载电阻的阻值使棒中的电流强度 $I$ 保持恒定。导体棒一直在磁场中运动。若不计导轨电阻，求此过程中导体棒上感应电动势的平均值和负载电阻上消耗的平均功率。

\onepicture[w=4.5cm, align=right]{figs/24.png}

%% answer:
\jdanswer{
\begin{enumerate}
	\item
	$\overline{E}=\frac{1}{2}Bl(v_{0}+v_{1})$
	\item
	$\overline{P}_{2}=\frac{1}{2}BIl(v_{0}+v_{1})-I^{2}r$
\end{enumerate}
}

%% memo:
\memoanswer{
导体棒所受的安培力为 $F=IlB$，该力大小不变，棒做匀减速运动，因此在棒的速度由 $v_{0}$ 减小至 $v_{1}$ 的过程中，平均速度为 $\overline{v}=\frac{v_{0}+v_{1}}{2}$。

当棒的速度为 $v$ 时，感应电动势的大小为 $E=Blv$，棒中的平均感应电动势为 $\overline{E}=Bl\overline{v}=\frac{1}{2}Bl(v_{0}+v_{1})$。

导体棒中消耗的热功率为 $P_{1}=I^{2}r$，负载电阻上消耗的平均功率为 $\overline{P}_{2}=\overline{E}I-P_{1}=\frac{1}{2}BIl(v_{0}+v_{1})-I^{2}r$。
}

\item
%% number: 25
%% paperName: 2008年全国理综Ⅱ第 25 题 20 分
%% typeId: 6
%% type: 计算题
%% chapter: 曲线运动
%% point: 人造地球卫星
%% method: 无
%% score: 20
%% degree: 450
%% duplicateId: 0
%% body:
我国发射的“嫦娥一号”探月卫星沿近似于圆形的轨道绕月飞行。为了获得月球表面全貌的信息，让卫星轨道平面缓慢变化。卫星将获得的信息持续用微波信号发回地球。设地球和月球的质量分别为 $M$ 和 $m$，地球和月球的半径分别为 $R$ 和 $R_{1}$，月球绕地球的轨道半径和卫星绕月球的轨道半径分别为 $r$ 和 $r_{1}$，月球绕地球转动的周期为 $T$。假定在卫星绕月运行的一个周期内卫星轨道平面与地月连心线共面，求在该周期内卫星发射的微波信号因月球遮挡而不能到达地球的时间（用 $M$、$m$、$R$、$R_{1}$、$r$、$r_{1}$ 和 $T$ 表示，忽略月球绕地球转动对遮挡时间的影响）。

%% answer:
\jdanswer{
$t=\frac{T}{\pi}\sqrt{\frac{Mr_{1}^{3}}{mr^{3}}}\left(\arccos\frac{R-R_{1}}{r}-\arccos\frac{R_{1}}{r_{1}}\right)$
}

%% memo:
\memoanswer{
如图，$O$ 和 $O'$ 分别表示地球和月球的中心。在卫星轨道平面上，$A$ 是地月连心线 $OO'$ 与地月球面的公切线 $ACD$ 的交点，$D$、$C$ 和 $B$ 分别是该公切线与地球表面、月球表面和卫星圆轨道的交点。根据对称性，过 $A$ 点在另一侧作地月球面的公切线，交卫星轨道于 $E$ 点。卫星在 $BE$ 弧运动时发出的信号被遮挡。

\onepicture[w=9.0cm]{figs/25_an.jpg}

设探月卫星的质量为 $m_{0}$，万有引力常量为 $G$，根据万有引力定律有 $G\frac{Mm}{r^{2}}=m\left(\frac{2\pi}{T}\right)^{2}r$，$G\frac{mm_{0}}{r_{1}^{2}}=m_{0}\left(\frac{2\pi}{T_{1}}\right)^{2}r_{1}$，式中 $T_{1}$ 是探月卫星绕月球转动的周期。两式相比得 $\left(\frac{T_{1}}{T}\right)^{2}=\frac{M}{m}\left(\frac{r_{1}}{r}\right)^{3}$。

设卫星的微波信号被遮挡的时间为 $t$，则由于卫星绕月做匀速圆周运动，应有 $\frac{t}{T_{1}}=\frac{\alpha-\beta}{\pi}$，式中 $\alpha=\angle CO'A$，$\beta=\angle CO'B'$。由几何关系得 $r\cos\alpha=R-R_{1}$，$r_{1}\cos\beta=R_{1}$。

联立解得 $t=\frac{T}{\pi}\sqrt{\frac{Mr_{1}^{3}}{mr^{3}}}\left(\arccos\frac{R-R_{1}}{r}-\arccos\frac{R_{1}}{r_{1}}\right)$。
}

\end{enumerate}

\end{document}
